\section{Conclusion}

We formulated visual transition credit assignment as a central problem in procedural skill evolution for partially observable embodied agents. A failed trajectory does not by itself identify whether a reusable skill is incorrect; the responsible source may instead be the current visual belief, persistent action knowledge, or a frozen executor that failed to follow a valid procedure. VISTA-Skill addresses this ambiguity by expressing visual belief, action effects, and procedural skills in a shared graph-predicate space and by comparing skill-predicted effects with an information-isolated, evidence-supported belief transition. Its hierarchical attribution mechanism first selects the appropriate update target and only then performs evidence-bound field revision and paired candidate verification. Experiments on EB-HAB and EB-NAV show that VISTA-Skill improves \textcolor{blue}{[TASK RESULTS]}, increases beneficial-update precision by \textcolor{blue}{[X]}, and reduces harmful-update rate by \textcolor{blue}{[Y]} over the open-source EmbodiSkill baseline and matched trajectory-level evolution controls. These results suggest that reliable embodied self-evolution depends not only on how a skill is optimized, but also on whether the available evidence warrants optimizing the skill at all.
